\section{Further extensions of studied dynamics}\label{sec:further_extensions_of_studied_dynamic}
The dynamics studied in this thesis were introduced to explain gaze, but each of their components is a general mechanism for exploring a sparse and unknown landscape. Nonlocal superdiffusive jumps give access to distant regions, momentum-driven local oscillations decorrelate samples within a region, a history-dependent kernel records the information gathered and discourages revisiting, and theta-rhythmic transitions organise when the sampler moves between regions of interest. Nothing in their construction is specific to vision, and the model can be extended to a state space of any dimension. In this section we outline how these components might carry over to two settings in machine learning and artificial intelligence: sampling from difficult probability distributions, and pursuing goals that change in time. These are proposals for future work and are not tested here.

\subsection{Dynamical goals}\label{subsec:dynamical_goals}
In the search tasks considered in this study, the goal is fixed: the landscape $V_0(\mathbf{x})$ and the target regions are set before the search begins, and only the observer's belief about it changes. Generally---including in machine learning---goals will be changing in time, as targets move, rewards shift, or agents switch tasks. A searcher must then remember where it has looked, to avoid wasting samples, and account for fact that old information may be less useful than new information. This motivates time-dependent landscapes $V_0(\mathbf{x}, t)$ which are intrinsically dynamical, beyond changing beliefs about the landscape. The information kernel already provides memory, and so including a mechanism of forgetting is simple; let the kernel effect decay back down to zero over time. Writing the potential as the current goal plus a kernel term recording the search history, $V_t(\mathbf{x}) = V_0(\mathbf{x}, t) + K_t(\mathbf{x})$, we let the kernel effect decay at a rate $\lambda$:
\begin{equation}
	\partial_t K_t(\mathbf{x}) = \frac{1}{\tau}k(\mathbf{x} - \mathbf{x}_t) - \lambda K_t(\mathbf{x}), \qquad K_0(\mathbf{x}) = 0.
\end{equation}
The case $\lambda = 0$ recovers the present model. For $\lambda > 0$, kernel deposition decays exponentially with memory rate $1/\lambda$ in units of time. This introduces two parameters into kernel dynamics, one ($\tau$) which defines the theta-rhythmic rate of information gathering, and the second ($\lambda$) which defines the decay rate of this information in dynamical tasks. This model can be used to make rigorous predictions about the effect of working memory on theta-scales in behavioural experiments involving dynamical goals, such as the classic experiment requiring tracking a moving pendulum \cite{law_eye-tracking_1973}.


\subsection{Superdiffusive probabilistic sampling}\label{subsec:superdiffusive_probability_sampling}
A central task in Bayesian inference and generative modelling is to draw samples from a distribution $\pi(\mathbf{x}) \propto \exp[-V(\mathbf{x})]$ that may be multimodal, with well-separated modes, or heavy-tailed \cite{latuszynski_mcmc_2025, betancourt_conceptual_2018}. FNS was developed in order to solve this problem in cortical activity dynamics, as exploring well-separated heavy-tailed distributions are the central dynamics of neural sampling (see \cref{subsec:neur_sampling} \cite{qi_fractional_2022}). From the field of abstract probabilistic sampling, standard gradient-based Markov chain Monte Carlo (MCMC) methods are driven by Brownian noise and therefore explore locally. A well-understood short-coming of these methods is that the time to escape a mode grows exponentially with the height of the surrounding barrier, so samplers can remain trapped in a single mode for very long periods if not tuned correctly \cite{latuszynski_mcmc_2025}. Furthermore, heavy-tailed distributions give small gradients, which slow the exploration, and often fail to constrain standard MCMC algorithms \cite{behme_levy_2026}.

Each component of the model addresses part of this problem, and this chapter has shown that each component efficiently provides an efficient strategy in searching through a complex multimodal landscape. First, driving the dynamics with $\alpha$--stable Lévy noise ($\alpha < 2$) replaces local normally diffusive exploration with global superdiffusion, which increases the access to far displaced modes of interest. For the Lévy flight of this model, escape from a mode is governed mainly by the distance to the next mode rather than by the height of the barrier between them \cite{behme_levy_2026, ijcai2018p419}. Jumps between separated modes occur at a rate that does not decay exponentially with the barrier, but polynomially with the separation \cite{behme_levy_2026, ijcai2018p419}. Second, the momentum coupling produces fast local oscillations within a mode. As in Hamiltonian Monte Carlo, these decorrelate nearby samples, so that a fixed number of samples covers more of the local structure of $\pi$ \cite{ijcai2018p419, betancourt_conceptual_2018}. Third, the dynamical-information kernel of \cref{sec:dynamical-information_adaptation} adds a history-dependent term to the potential, which penalises revisited regions. For extended sampling, this kernel can be allowed to relax down after the mode has been explored. This kernel deposition closely resembles metadynamics \cite{tiwary_time-independent_2015}, in which deposited bias fills explored minima and pushes the sampler towards new ones. Finally, the regularity of a sampler with nested timescales has not been presented in the literature so-far, and this would represent a novel addition, allowing a decoupling of local sampling and the global exploration which modulates it, rather than through a hand-tuned mixture of local and global proposals, or a single local sampler which fails to explore multiple modes.

All of these dynamics can naturally produce an efficient, minimally tuned algorithm, only requiring a spatial scale defined by the covariance of an explored target, from which to set noise strength, and ensure local harmonic exploration. Developing a superdiffusive sampler is a active area of research in modern MCMC literature \cite{behme_levy_2026, simsekli_fractional_nodate, ijcai2018p419}, and constructing a model based on the present equations of motion (\cref{eq:fns_system}) is a valuable addition to this literature. A full study in this direction would need to consider the following two components of the model: firstly, the information kernel adjusts the sampled distribution, and this would need be accounted for through, for instance, reweighting as in metadynamics \cite{tiwary_time-independent_2015, bussi_using_2020, barducci_well-tempered_2008}, and secondly, calculation of the full fractional derivative in the gradient drift term may be computationally prohibitive. There exist creative methods to account for both considerations, but such a study is left to further work.